%% When Consensus Lies — CSCW (PACM HCI) full paper
%% Format: ACM acmart, acmsmall (PACM HCI journal single-column). `review,anonymous` for submission.
%% Build: pdflatex -> bibtex -> pdflatex x2  (bib: references.bib)
%%
%% ⚠️ CSCW 2026 LLM-writing policy: this file is a Manager-prepared SCAFFOLD. The author MUST rewrite the
%% prose in their own voice, INDEPENDENTLY VERIFY EVERY CITATION (some bib entries are approximate), and
%% follow the venue's disclosure policy. Do not submit as-is. See paper/research/2026-07-20-cscw-writing-style-guide.md.
\documentclass[acmsmall,screen,review,anonymous]{acmart}

\acmJournal{PACMHCI}
\acmVolume{9}
\acmNumber{CSCW}
\acmArticle{XXX}
\acmMonth{4}
\acmYear{2026}
\acmDOI{10.1145/XXXXXXX}
\setcopyright{acmlicensed}

\usepackage{booktabs}
\usepackage{amsmath}
\usepackage{array}
\usepackage{xcolor}
\newcommand{\todo}[1]{\textcolor{red}{[\textbf{TODO:} #1]}}
\newcommand{\cd}{\textsf{CD}}
\newcommand{\iperp}{$I_\perp$}
\newcommand{\neff}{$n_{\mathrm{eff}}$}

\begin{document}

\title{When Consensus Lies: Fake Redundancy and Silent Convergent Delusion in AI-Assisted Collective Judgment}

\author{Anonymous Author(s)}
\affiliation{%
  \institution{Anonymous Institution}
  \city{}
  \country{}
}

\begin{abstract}
Groups and organizations increasingly treat \emph{agreement among multiple AI systems} as corroboration:
we sample a language model repeatedly, pool several models from different vendors, or stage multi-agent
``debates,'' and read the resulting consensus as a warrant for a decision. The epistemic value of such
consensus rests on an assumption inherited from the wisdom of crowds---that the judgments being aggregated
carry \emph{sufficiently independent} evidence. We show that this assumption can fail silently. In a
pre-registered study on a reconstructed, executable-gold benchmark of \emph{underspecified} tasks, we
manipulate a single factor---whether the decisive disambiguating information lies \emph{outside} or
\emph{inside} the retained prompt---and measure not just accuracy but the \emph{correlation} of errors
across agents. When the disambiguator is outside the prompt, independent agents across capability tiers and
across model families (OpenAI, Anthropic, Google) converge on the \emph{same wrong} interpretation
(convergent-delusion rate $\approx\!0.53$), almost never signaling that anything is missing (abstention
$\approx\!0.02\%$); a nominal $k$-agent ensemble supplies the evidence of $\approx\!1$ independent judgment
($\neff\!=\!1.10$). When the same information is inside the prompt, every tier---including legacy-weak
models---resolves it (rate $=\!0.00$). A second benchmark constructed by a \emph{different} model family
rules out a shared-prior-with-the-constructor artifact. We name this failure \emph{fake redundancy}: under a
shared information boundary, nominally diverse agents form an \emph{interpretive monoculture}, and consensus
stops being evidence of correctness. We derive implications for how collaborative systems should aggregate,
display, and place trust in AI judgments---measuring \emph{dependence} rather than counting \emph{agreement}.
\end{abstract}

\begin{CCSXML}
<ccs2012>
 <concept><concept_id>10003120.10003130.10011762</concept_id>
  <concept_desc>Human-centered computing~Collaborative and social computing theory, concepts and paradigms</concept_desc>
  <concept_significance>500</concept_significance></concept>
 <concept><concept_id>10003120.10003130.10011765</concept_id>
  <concept_desc>Human-centered computing~Collaborative and social computing systems and tools</concept_desc>
  <concept_significance>300</concept_significance></concept>
 <concept><concept_id>10010147.10010178</concept_id>
  <concept_desc>Computing methodologies~Artificial intelligence</concept_desc>
  <concept_significance>300</concept_significance></concept>
</ccs2012>
\end{CCSXML}
\ccsdesc[500]{Human-centered computing~Collaborative and social computing theory, concepts and paradigms}
\ccsdesc[300]{Human-centered computing~Collaborative and social computing systems and tools}
\ccsdesc[300]{Computing methodologies~Artificial intelligence}

\keywords{collective intelligence, wisdom of crowds, consensus, common ground, algorithmic monoculture,
human--AI collaboration, large language models, multi-agent systems, error correlation, automation bias,
pre-registration}

\maketitle

%% =====================================================================
\section{Introduction}
\label{sec:intro}
When people must decide under uncertainty, they seek \emph{corroboration}: they consult multiple sources,
solicit second opinions, and trust a conclusion more when independent parties agree. This practice is
increasingly automated. Software teams triage incidents with several model ``reviewers''; analysts pool
outputs from different vendors' assistants; emerging agentic pipelines sample a model many times and take
the majority, or stage multi-agent debate, and surface the result to a person as ``$5/5$ agents agree.''
In each case, \emph{agreement among AI systems is read as a signal that the answer can be trusted}.

The epistemic warrant for this practice is old and specific. From the Condorcet Jury Theorem to the wisdom
of crowds, aggregation improves reliability \emph{only when} the judgments being combined are competent and
carry \emph{sufficiently independent} errors~\cite{condorcet,ladha1992}; social influence that erodes that
independence produces confident convergence on error rather than accuracy~\cite{lorenz2011}. Aggregating AI
judgments imports this warrant, but it also imports its precondition---and it is far from obvious that
querying a model five times, or pooling five vendors, produces five \emph{independent} interpretations of
the task.

We show that it often does not, and we identify \emph{when}. Our study concerns \emph{underspecified}
inputs: inputs from which a decisive disambiguating clause has been omitted, so more than one
interpretation is consistent with what remains---an everyday condition when context is lost across a handoff
between people, tools, and agents. We find a sharp, pre-registered, two-regime phenomenon governed not by
how capable the models are but by \emph{where the missing information lives}. When the disambiguator lies
\emph{outside} the retained prompt, independent agents---across capability tiers, and across three model
families---\emph{silently} converge on the \emph{same wrong} interpretation. They do not make independent
errors that average away; they make the \emph{same} error, in unison and with confidence, and a nominal
$k$-agent ensemble collapses to roughly one effective independent judgment. When the same disambiguator is
present \emph{inside} the prompt, every tier resolves it. We call the wrong-consensus failure
\emph{convergent delusion}, and the structural condition that produces it \emph{fake redundancy}: apparent
redundancy that supplies no independent evidence because all agents inherit the same lossy representation of
the task.

This reframes a question at the heart of CSCW. Aggregating AI advice is not merely an accuracy problem; it
is a question about whether our collaborative infrastructures preserve the \emph{epistemic independence}
that makes consensus meaningful. When nominally diverse agents share an information boundary, they can form
an \emph{interpretive monoculture}~\cite{kleinberg2021monoculture} whose agreement is an artifact of a shared
prior rather than corroborating evidence---and, because the failure is silent, nothing in the system's
behavior distinguishes it from genuine agreement.

\paragraph{Contributions.} This paper makes four contributions.
\begin{itemize}
  \item \textbf{Empirical.} Pre-registered evidence that independently queried LLM agents, spanning
  capability tiers \emph{and} three model families, silently converge on the same incorrect interpretation
  when a decisive disambiguator is omitted from the retained prompt---and that this failure vanishes when
  the disambiguator is present. The failure is capability-invariant and cross-family, not an artifact of a
  single model or vendor.
  \item \textbf{Methodological.} An \emph{executable-gold} protocol that manipulates whether the
  disambiguator crosses the prompt boundary and measures \emph{correlated error}---not merely individual
  accuracy---so that correctness is established independently of model agreement and no LLM judge is in the
  loop.
  \item \textbf{Conceptual.} \emph{Fake redundancy} as a failure mode of AI aggregation, operationalized
  through classical dependence estimators (pairwise wrong-answer agreement, intraclass correlation, and
  effective ensemble size), showing when a nominally $k$-agent system supplies $\approx\!1$ independent
  judgment.
  \item \textbf{Design.} Implications for AI-mediated collaborative decision systems: consensus displays
  should expose interpretive and evidential \emph{dependence} rather than agent count; systems should elicit
  competing interpretations before aggregating answers, preserve contextual provenance across handoffs, and
  route to human review when agreement is high but contextual completeness is uncertain.
\end{itemize}
We report our pre-registered hypotheses honestly, including those that were not supported, and treat the
capability/family nulls as theoretically informative: they are what makes the failure a property of the
\emph{information boundary} rather than of any model.

%% =====================================================================
\section{Consensus, Grounding, and AI-Mediated Collective Judgment}
\label{sec:related}
We organize related work around the assumptions our study tests, rather than around technologies.

\subsection{When does aggregation help? Independence and its erosion}
The Condorcet Jury Theorem's optimistic conclusion---that majority accuracy approaches one as a group
grows---depends on competence \emph{and} sufficiently independent votes; correlation among votes changes,
and can invert, the result~\cite{condorcet,ladha1992,dietrich2013}. Ensemble learning reaches the same
conclusion formally: ensemble error is average member error minus a diversity term, and a common-error
floor persists as members correlate~\cite{krogh1994,tumer1996,breiman2001}. For $n$ exchangeable errors
with pairwise correlation $\rho$, the variance of the mean is $\sigma^2(\rho + (1-\rho)/n)$: as
$\rho\!\to\!1$, $n$ nominal contributors provide one effective independent signal. Empirically, social
influence can reduce the diversity that crowds rely on \emph{without} improving accuracy, yielding confident
convergence around error~\cite{lorenz2011}. \emph{Unresolved premise our study tests:} whether independently
\emph{queried} AI agents form independent \emph{interpretations} when all inherit the same underspecified
representation.

\subsection{Grounding and information boundaries in cooperative work}
Grounding is the process by which collaborators establish mutual understanding sufficient for present
purposes~\cite{clark1991grounding}. We treat the retained prompt as the agents' shared information
environment and the deleted disambiguator as a \emph{grounding-resource boundary}: information available
in-prompt enters the interpretive process, information lost at the handoff does not, and downstream
agreement cannot repair an upstream grounding failure. We do not claim that non-interacting agents establish
common ground with one another; we claim that their \emph{common information environment} determines whether
convergence is warranted. \emph{Unresolved premise:} whether corroboration across agents can substitute for
the grounding lost at a context handoff.

\subsection{Reliance on aggregated AI advice}
Overreliance on automation is a long-standing concern~\cite{parasuraman1997,skitka1999}; explanations help
mainly when they lower the cost of \emph{verifying} the machine, not when they merely assert
confidence~\cite{vasconcelos2023explanations}. A ``$5/5$ agents agree'' display is precisely the kind of cue
that can increase reliance while adding little independent evidence. \emph{Unresolved premise:} whether
consensus displays convey evidence or only the \emph{appearance} of it.

\subsection{Algorithmic monoculture and shared priors}
When many decision-makers rely on the same model or data, individually accurate systems can reduce
system-level welfare through correlated outcomes~\cite{kleinberg2021monoculture}. We identify a related
\emph{within-ensemble} phenomenon: even nominally heterogeneous models can form an interpretive monoculture
when decisive context is uniformly absent, and recent evidence shows LLM errors are correlated across
providers and architectures, especially among stronger models~\cite{kim2025correlated}, with sycophancy and
shared preference training as plausible sources~\cite{perez2022,sharma2024}.

\subsection{LLM ensembles and multi-agent systems}
Self-consistency samples multiple chains and votes~\cite{wang2023selfconsistency}; multi-agent debate can
help~\cite{du2023debate} but also converge to shared misconceptions~\cite{estornell2024,liang2024}, conform
to majorities~\cite{choi2025conformity}, and degrade under peer pressure~\cite{wynn2025talk}. Underspecified
prompts are common and destabilizing~\cite{yang2025underspecification}. Our contribution is not to
re-demonstrate that consensus can be wrong; it is to \emph{isolate a controlled axis}---the location of the
disambiguator---that switches the failure on and off while holding the task, answer space, and verifier
fixed, and to connect it to the CSCW question of when aggregated AI judgment can be trusted.

%% =====================================================================
\section{Research Questions and Hypotheses}
\label{sec:rqs}
We pre-registered the following (\S\ref{sec:study}; deviations in Appendix~\ref{app:deviations}).
\textbf{RQ1 (regime):} does convergent delusion depend on the \emph{location} of the disambiguator?
We predict high wrong-consensus when it is external (\textbf{H1\_external}) and near-zero when it is
in-prompt (\textbf{H2\_derivable}). \textbf{RQ2 (redundancy):} does adding agents add independent evidence?
We test aggregation vs.\ a single agent (\textbf{H1a}) and quantify effective ensemble size.
\textbf{RQ3 (capability/family):} is the failure explained by weak models or a single vendor? We compare
capability tiers and model families. \textbf{RQ4 (dose):} does wrong-consensus increase with the number of
omitted clauses $k$ (\textbf{H1b})? We include a fully specified $k{=}0$ control (\textbf{R1a}), an
interpretation-uniform null (\textbf{R1b}), and a cross-family \emph{construction} control (\textbf{R2}).

%% =====================================================================
\section{Executable-Gold Benchmark}
\label{sec:benchmark}
\subsection{Construct and design}
A task is \emph{underspecified} when a decisive clause is removed so that multiple interpretations remain
consistent with the retained input. Each task fixes a target interpretation $I_0$ and enumerated
alternatives (foils) $I_1,\dots$; an answer is labeled deterministically by \emph{which} interpretation it
realizes---via an executable \texttt{check()}---or \iperp{} if it is off-axis or unparseable. Critically,
correctness and the interpretation label are established \emph{independently of model agreement}: no LLM
judges any output.

We reconstruct 54 task variants across three domains (code specification, data analysis, policy
question-answering). Items are \emph{reversed} so the natural default an unaware solver would produce is a
\emph{foil}, and each foil differs from $I_0$ on exactly one convention axis (e.g., calendar vs.\ fiscal
quarter; floor vs.\ truncating integer division). The set contains 40 \textbf{H1\_external} items (the
disambiguator requires outside knowledge) and 14 \textbf{H2\_derivable} items (the disambiguator is present
in the retained prompt), each at ambiguity level $k\in\{0,1,2,3\}$, where $k{=}0$ is fully specified.
\todo{Table: domains $\times$ axes $\times$ example; unique-$I_0$ fairness; construction validation.}

\subsection{Construct validity}
An external disambiguator must be realistic omitted context, not adversarial trivia, and $I_0$ must be
uniquely established by an out-of-prompt convention. We validated each item with an empirical default check
(capable models, unaware, produce the foil) and, for the regime tag, that in-prompt disambiguators are
resolved. \todo{report the default-check summary.}

%% =====================================================================
\section{Pre-registered Study}
\label{sec:study}
\subsection{Models, tiers, and families}
We evaluate a frontier roster spanning three families---OpenAI, Anthropic, Google---grouped into model
classes: a homogeneous single-model baseline, a heterogeneous cross-family pool, a reasoning tier, and a
legacy-weak tier. \emph{Independence protocol:} agents are queried independently (no inter-agent messages in
the primary conditions) so that any error correlation reflects shared priors and shared inputs, not
conversational influence.

\subsection{Conditions and estimands}
Primary aggregation conditions: single-agent baseline; self-consistency ($k{=}5$); homogeneous
majority/debate ($N{=}5$); and heterogeneous cross-family debate. A scheduled secondary sensitivity pass adds
self-consistency $k{=}10$ \todo{(running; \S\ref{sec:results})} and a verifier condition. Our primary
outcome is \emph{convergent delusion} \cd: the share of agents on the single modal \emph{wrong enumerated}
interpretation (\iperp{} stays in the denominator but cannot be the convergent label), with two
pre-registered sensitivity variants (\iperp{}-eligible; \iperp{}-dropped). We additionally report dependence
estimators over each cell's agent labels---pairwise wrong-answer agreement, Fleiss' $\kappa$, the intraclass
correlation of the wrong-answer indicator, and the effective ensemble size $\neff=n/(1+(n-1)\bar\rho)$---and
an independence-calibrated counterfactual that resamples each agent from its own marginal accuracy. A
rule-based (non-LLM) detector flags whether an agent abstained or requested clarification. Each cell is
replicated over three collision-free seeds.

\subsection{Analysis and open science}
Hypotheses, the primary metric, the regime criterion, decision rules, and robustness nulls were
git-committed and frozen before the confirmatory run; every subsequent change is a timestamped, signed
amendment reported in Appendix~\ref{app:deviations}. Provenance was separated by construction: analysis
artifacts were implemented and \emph{independently, adversarially audited by different model families}, and
all labels derive from deterministic executable gold. \todo{cite prereg + amendments + release plan.}

%% =====================================================================
\section{Results}
\label{sec:results}
The confirmatory run comprised 5{,}832 aggregation jobs (16{,}461 agent executions) over 54 items, three
seeds, and the primary conditions. We report answer-first findings.

\subsection{F1: Omitting external context produces accurate-looking but systematically wrong consensus}
On \textbf{H1\_external}, \cd{} $=0.53$ (0.60 among parseable answers): a majority of independent agents land
on the \emph{same wrong} interpretation. On \textbf{H2\_derivable}, \cd{} $=0.00$: with the disambiguator in
the prompt, no parseable agent converges on a wrong enumerated interpretation. The regime gap is
$\Delta\cd=0.53$, consistent across all three metric variants (Table~\ref{tab:regime}). A within-design
manipulation check makes this decisive: on fully specified $k{=}0$ items \cd{} $=0.00$, and on $k\ge1$ items
\cd{} $=0.82$ (difference 95\% CI $[0.73,0.90]$)---the same pipeline, labeler, and metric yield zero when
the prompt is complete and high when a single clause is deleted (\textbf{R1a}), ruling out a labeler or
foil-count artifact.

\begin{table}[t]
\centering
\caption{Convergent delusion by regime and model class (primary conditions). \cd{} is capability-invariant
on H1\_external and zero on H2\_derivable.}
\label{tab:regime}
\begin{tabular}{llccc}
\toprule
Regime & Model class & \cd{} (primary) & \cd{} (parseable) & \iperp{} rate \\
\midrule
H1\_external & homogeneous   & 0.536 & 0.562 & 0.112 \\
H1\_external & heterogeneous & 0.544 & 0.558 & 0.106 \\
H1\_external & reasoning     & 0.539 & 0.625 & 0.107 \\
H1\_external & weak          & 0.517 & 0.613 & 0.129 \\
\midrule
H2\_derivable & homogeneous  & 0.000 & 0.000 & 0.152 \\
H2\_derivable & heterogeneous& 0.000 & 0.000 & 0.048 \\
H2\_derivable & reasoning    & 0.000 & 0.000 & 0.016 \\
H2\_derivable & weak         & 0.000 & 0.000 & 0.067 \\
\bottomrule
\end{tabular}
\end{table}

\subsection{F2: Capability and model-family diversity do not restore interpretive independence}
The H1\_external failure is essentially flat across model classes (Table~\ref{tab:regime}): frontier
reasoners are trapped as reliably as legacy-weak models, and a heterogeneous cross-family pool converges as
strongly as a single model sampled repeatedly (heterogeneous vs.\ homogeneous $\Delta\cd=-0.003$, 95\% CI
$[-0.078,+0.058]$; not the ordering we predicted). Vendor and capability diversity did not translate into
\emph{viewpoint} diversity: the shared prior spans families.

\subsection{F3: Adding agents adds little effective redundancy (fake redundancy)}
On H1\_external, pairwise wrong-answer agreement is $0.98$, Fleiss' $\kappa=0.97$, and the intraclass
correlation of the wrong-answer indicator is $0.89$; the effective ensemble size is $\neff=1.10$. An
independence-calibrated counterfactual yields observed$-$independent $\cd=+0.245$ (95\% CI $[0.234,0.257]$):
observed convergence significantly exceeds what genuinely independent agents at the same accuracy would
produce. Consistently, the pre-registered aggregation-vs-single contrast (\textbf{H1a}) is inconclusive
(agreement neither amplifies nor removes \cd{}), as expected when the ensemble has already collapsed at
$k{=}5$. \emph{A nominal $k$-agent system here supplies about one independent judgment.} \todo{SC $k{=}10$
sensitivity (running): confirm $\neff\!\to\!1$ persists at higher $k$.}

\subsection{F4: The failure is silent; retaining the disambiguator eliminates it}
On H1\_external the abstention / clarification-request rate is $\approx0.02\%$ while \cd{} $=0.53$: agents
are confidently wrong and almost never signal missing information, so consensus carries no internal warning.
Wrong-consensus also rises with the number of omitted clauses $k$ (\textbf{H1b} supported; $k$-coefficient
$0.31$, 95\% CI $[0.18,0.45]$). The constructive boundary condition is that \emph{presence} of the
disambiguator (H2) removes the failure entirely.
\todo{finalize abstention once the rule-based detector's human-validation sample is checked.}

\subsection{F5: The convergence is not an artifact of the benchmark's constructor}
A skeptic might argue that convergence reflects a prior shared with whoever \emph{built} the traps. We
reconstructed a held-out subset of H1 items with a \emph{different} constructor family (Anthropic) and
re-ran it: models from the \emph{other} families (OpenAI, Google) converge on these Anthropic-constructed
foils at single-agent \cd{} $=0.72$ (0.87 on construct-valid items). The convergence persists under
cross-family construction (\textbf{R2}); the pre-registered R2 aggregation-amplification metric was, like
R1b below, statistically inconclusive on the primary metric, so we rest the confound rebuttal on the direct
cross-family single-agent persistence and report the metric honestly.

\subsection{Robustness and honestly reported nulls}
The within-design $k{=}0$ control (R1a) is strongly supported (above). An interpretation-uniform null (R1b)
is inconclusive on the primary metric---binary items give a high random baseline---but supported among
parseable answers; the excess-convergence claim rests on R1a and the independence-calibrated counterfactual,
which are both decisive. The H2 reasoner$\times$capability \emph{interaction} was not estimable (the crossed
random-effects model did not converge on the full interaction) and was, as anticipated, an expected null;
the well-powered regime \emph{main} effect carries the claim. \todo{note the two items with incomplete foil
sets, handled by the drop-\iperp{} sensitivity.}

%% =====================================================================
\section{Discussion}
\label{sec:discussion}
\subsection{Consensus without epistemic independence}
Our results identify a boundary condition for applying collective-intelligence assumptions to AI ensembles.
Aggregation earns its epistemic warrant from independence; under a shared information boundary, that warrant
is absent even when the aggregated systems are nominally diverse. The failure is not that models are weak
(reasoners fail identically), nor that they are homogeneous (cross-family pools fail identically), but that
they share the \emph{same lossy representation of the task}. Consensus then measures the concentration of a
shared prior, not the corroboration of independent evidence.

\subsection{Shared-input failure, not groupthink}
The outcome resembles groupthink---unanimous, confident error~\cite{janis1972}---but the mechanism differs.
Classical groupthink arises from cohesion, conformity pressure, and suppressed dissent among interacting
members; our agents do not communicate in the primary conditions, yet still converge. This points to a
distinct pathway---a \emph{shared-input / representational monoculture}---in which unanimity requires no
social influence at all. Naming this precisely matters: interventions for groupthink (devil's advocacy,
structured dissent \emph{among} members) will not address a failure that occurs before any interaction.

\subsection{Information loss as a sociotechnical failure}
Because the decisive information is lost at a handoff (a grounding-resource boundary), downstream
aggregation cannot recover it. The problem is therefore located in the \emph{pipeline}, not the model: it is
a property of how collaborative systems preserve (or discard) context and provenance across people, tools,
and agents. The silence of the failure compounds this: without an internal warning, more agreement reads as
more confidence.

%% =====================================================================
\section{Implications for CSCW and Collaborative AI Systems}
\label{sec:implications}
We separate recommendations by evidential strength. Directly supported by our findings: (1)~\emph{Measure
dependence, not agreement.} Systems that aggregate AI judgments should estimate error dependence during
validation and surface effective ensemble size, not agent count; a ``$5/5$ agree'' display can be
\emph{replaced} by a dependence-aware uncertainty display. (2)~\emph{Aggregate interpretations before
answers.} Because the failure is interpretive, eliciting each agent's task \emph{interpretation} before
soliciting an answer can expose latent forks that a final vote hides. (3)~\emph{Route on contextual
completeness.} Trigger human review when consensus is high \emph{but} contextual completeness is
uncertain---high agreement is least informative exactly when the input is underspecified. Supported by our
findings together with prior CSCW work on reliance~\cite{vasconcelos2023explanations}: (4)~\emph{Diversify
information, not just vendors.} Viewpoint diversity requires diverse \emph{inputs}, retrieval paths, and
representations, not merely different model providers. (5)~\emph{Preserve provenance across handoffs}, so
that the context needed to disambiguate is not silently dropped. (6)~\emph{Elicit clarification.} Since
agents rarely ask, systems---or their designers---must build the clarification-seeking that the models omit.
(7)~\emph{Maintain dissent channels} rather than forcing early convergence. (8)~\emph{Audit for interpretive
monoculture} across a fleet, not only per-model accuracy.

%% =====================================================================
\section{Limitations, Ethics, and Open Science}
\label{sec:limitations}
\paragraph{Limitations.} Our tasks are single-episode and non-interactive; real collaborative pipelines
involve iteration and human participation, which we do not test. Claims about human trust, teams, and
high-stakes domains are \emph{implications}, not tested deployments. Construct validity rests on the
reversed-benchmark design and the default check; two items had incomplete foil sets and are handled by the
pre-registered drop-\iperp{} sensitivity. The abstention detector is rule-based and pending human
validation. Generalization beyond the three domains and the tested roster is future work.
\paragraph{Ethics.} \todo{Benchmark construction; API/provider terms; compute/environmental cost (runs were
\$0 on a local proxy); release risks; no human subjects.} \paragraph{Open science.} We pre-registered the
study and will release the benchmark, executable evaluators, prompts, model/version records, analysis code,
and the amendment ledger. \todo{link on acceptance.}

%% =====================================================================
\section{Conclusion}
\label{sec:conclusion}
When collaborative systems aggregate AI judgments, they inherit an assumption that is easy to state and easy
to violate: that agreement reflects independent evidence. Under underspecification, that assumption fails
silently---nominally diverse agents form an interpretive monoculture, a $k$-agent ensemble collapses to one
effective voice, and confident consensus becomes an artifact of a shared prior rather than a warrant for a
decision. The fix is not more agents but a different question. For CSCW, designing trustworthy AI-mediated
collective judgment means measuring \emph{independence}, preserving \emph{context}, and treating consensus
as a hypothesis to be checked---not as proof.

\begin{acks}
\todo{Anonymized for review.}
\end{acks}

\bibliographystyle{ACM-Reference-Format}
\bibliography{references}

\appendix
\section{Deviations from Pre-registration}
\label{app:deviations}
\todo{Import the deviations ledger (paper/writing/2026-07-18-deviations-from-preregistration.md): frozen
prereg + amendments A01--A11, and the honestly reported H1a / R1b / heterogeneity / H2-interaction outcomes.}

\end{document}
